\chapter{DRI Concepts}
\label{ch:dri-concepts}

% ─────────────────────────────────────────────────────────────────────────────
\section{The DRI Model}
% ─────────────────────────────────────────────────────────────────────────────

The \textbf{Directly Responsible Individual} (DRI) model is a single-owner accountability pattern. At any moment during a runbook execution, exactly one person is the DRI: they have authority to make decisions, approve gates, and sign off on evidence. The DRI is accountable for the outcome---if a change fails, the DRI answers for it. If an incident escalates, the DRI owns the response.

\paragraph{Single Owner Invariant.}
A runbook execution has exactly one DRI at any given time. The DRI may change during the run (e.g., when an incident-commander takes over), but at no point are there zero DRIs or multiple conflicting DRIs.

\paragraph{DRI vs. Executor.}
The DRI is not always the person executing the steps. For example:
\begin{itemize}
  \item A database engineer may run the migration steps, but the DRI is the lead engineer who approved the plan.
  \item During an incident, multiple responders may be troubleshooting, but the incident-commander is the DRI.
\end{itemize}

The DRI is accountable, not necessarily hands-on-keyboard.

% ─────────────────────────────────────────────────────────────────────────────
\section{Role Taxonomy}
\label{sec:role-taxonomy}
% ─────────────────────────────────────────────────────────────────────────────

The \texttt{gert.ops} Kit defines six role types. Each role has specific permissions and responsibilities.

\subsection{dri — Directly Responsible Individual}

\textbf{Permissions:}
\begin{itemize}
  \item Execute all steps (unless a step requires a more specific role)
  \item Submit evidence
  \item Approve gates if listed as an approver
  \item Cancel the run
\end{itemize}

\textbf{Responsibilities:}
\begin{itemize}
  \item Own the outcome of the runbook execution
  \item Make decisions when the runbook requires judgment calls
  \item Sign off on the final evidence bundle
\end{itemize}

\textbf{Assignment:} The DRI is declared in \texttt{meta.roles.dri}. If no DRI is declared, the person starting the run becomes the implicit DRI.

\subsection{approver}

\textbf{Permissions:}
\begin{itemize}
  \item Sign off on \texttt{ops.approval} gates
  \item View the run status and evidence
\end{itemize}

\textbf{Responsibilities:}
\begin{itemize}
  \item Review the evidence before signing off
  \item Reject approval if conditions are not met
\end{itemize}

\textbf{Assignment:} Approvers are listed in \texttt{meta.roles.approver} (may be a person, team, or on-call alias). Specific approval gates may override this with a gate-specific approver list.

\subsection{change-manager}

\textbf{Permissions:}
\begin{itemize}
  \item Approve or reject the run before it starts
  \item View the final evidence bundle
\end{itemize}

\textbf{Responsibilities:}
\begin{itemize}
  \item Ensure the change request is properly authorized
  \item Verify that the run is scheduled during an approved change window
  \item Sign off on the change record after completion
\end{itemize}

\textbf{Assignment:} Declared in \texttt{meta.roles.change-manager}. If a \texttt{ops.change-request} step exists, the change-manager must approve before execution begins.

\subsection{incident-commander}

\textbf{Permissions:}
\begin{itemize}
  \item Take over as DRI when an incident is declared
  \item Assign tasks to responders
  \item Escalate or de-escalate the incident severity
\end{itemize}

\textbf{Responsibilities:}
\begin{itemize}
  \item Coordinate the incident response
  \item Decide when to execute rollback steps
  \item Declare the incident resolved
\end{itemize}

\textbf{Assignment:} Declared in \texttt{meta.roles.incident-commander}. When an \texttt{ops.incident} step is executed, the incident-commander becomes the active DRI for the duration of the incident.

\subsection{responder}

\textbf{Permissions:}
\begin{itemize}
  \item Execute steps that require the \texttt{responder} role
  \item Submit evidence
  \item Communicate with the incident-commander
\end{itemize}

\textbf{Responsibilities:}
\begin{itemize}
  \item Execute assigned triage or mitigation tasks
  \item Report status to the incident-commander
\end{itemize}

\textbf{Assignment:} Declared in \texttt{meta.roles.responder} (may be a team or on-call alias).

\subsection{observer}

\textbf{Permissions:}
\begin{itemize}
  \item View the run status and evidence in real-time
  \item Receive notifications when specific events occur
\end{itemize}

\textbf{Responsibilities:}
\begin{itemize}
  \item None. Observers are read-only participants.
\end{itemize}

\textbf{Assignment:} Declared in \texttt{meta.roles.observer}. Observers are typically stakeholders, managers, or compliance auditors who need visibility but not control.

% ─────────────────────────────────────────────────────────────────────────────
\section{Role Assignment in Runbooks}
% ─────────────────────────────────────────────────────────────────────────────

Roles are assigned in two places:

\paragraph{Runbook-level role declarations.}
The \texttt{meta.roles} block declares the default mapping of role names to identities:

\begin{minted}{yaml}
meta:
  roles:
    dri: alice@example.com
    approver: ops-team
    change-manager: change-board
    incident-commander: @on-call/sre
    responder: sre-team
    observer: [manager@example.com, compliance@example.com]
\end{minted}

\paragraph{Step-level role requirements.}
Individual steps may require a specific role to execute:

\begin{minted}{yaml}
steps:
  - id: migrate-database
    type: ops.cli
    requires_role: dri
    command: ./migrate.sh

  - id: verify-replication
    type: ops.manual
    requires_role: responder
    prompt: Check that all replicas are caught up
\end{minted}

If a step has \texttt{requires\_role}, only the person assigned to that role may execute the step. The runtime will block execution and emit an error if the current actor does not have the required role.

% ─────────────────────────────────────────────────────────────────────────────
\section{The Escalation Chain}
% ─────────────────────────────────────────────────────────────────────────────

If the primary DRI is unavailable (e.g., they go offline during an incident), responsibility must transfer to a backup. The \texttt{gert.ops} Kit supports escalation chains:

\begin{minted}{yaml}
meta:
  roles:
    dri: alice@example.com
  escalation:
    - primary: alice@example.com
      backup: bob@example.com
      timeout: 10m
    - primary: bob@example.com
      backup: @on-call/sre
      timeout: 5m
\end{minted}

\paragraph{How escalation works.}
\begin{enumerate}
  \item The runtime starts with the primary DRI (alice).
  \item If alice does not respond to an approval gate or manual step within 10 minutes, the runtime escalates to bob.
  \item If bob does not respond within 5 minutes, the runtime escalates to the on-call SRE.
  \item If the final backup does not respond, the run fails with a timeout error.
\end{enumerate}

\paragraph{Escalation triggers.}
Escalation is triggered by:
\begin{itemize}
  \item \texttt{ops.approval} timeout: the approval gate specifies a timeout, and no approver signs off before the deadline.
  \item \texttt{ops.manual} timeout: a manual step specifies a timeout, and the DRI does not mark it complete.
  \item SLA violation: the \texttt{meta.sla.max\_duration} is exceeded, and the run is still in progress.
\end{itemize}

% ─────────────────────────────────────────────────────────────────────────────
\section{SLA Concepts}
\label{sec:sla}
% ─────────────────────────────────────────────────────────────────────────────

SLA (Service Level Agreement) policies define the execution window and timeout behavior for a runbook.

\subsection{Execution Window}

The execution window is the time period during which the runbook is allowed to run. If the run is started outside the window, it is rejected:

\begin{minted}{yaml}
meta:
  sla:
    execution_window:
      start: "2025-01-15T02:00:00Z"
      end: "2025-01-15T06:00:00Z"
\end{minted}

This is useful for change management: a deployment runbook may only be executed during a scheduled maintenance window.

\subsection{Maximum Duration}

The maximum duration is the longest time the runbook is allowed to take:

\begin{minted}{yaml}
meta:
  sla:
    max_duration: 2h
\end{minted}

If the run exceeds this duration, the runtime triggers the escalation policy (either fail the run or escalate to a backup DRI).

\subsection{Escalation Policy}

The escalation policy defines what happens when an SLA timeout is reached:

\begin{minted}{yaml}
meta:
  sla:
    on_timeout: escalate  # or: fail
\end{minted}

\begin{description}
  \item[\texttt{fail}] — Terminate the run and emit a \texttt{run/failed} event.
  \item[\texttt{escalate}] — Transfer DRI responsibility to the next person in the escalation chain.
\end{description}

% ─────────────────────────────────────────────────────────────────────────────
\section{DRI Transfer}
% ─────────────────────────────────────────────────────────────────────────────

DRI transfer occurs when the primary DRI hands off responsibility to another person. This is common during incident response when shifts change or when an escalation timeout fires.

\paragraph{Explicit transfer.}
The DRI may explicitly transfer ownership:

\begin{verbatim}
$ gert exec transfer --run-id abc123 --to bob@example.com
\end{verbatim}

The runtime records the transfer in the trace:

\begin{minted}{json}
{
  "type": "governance/driTransfer",
  "timestamp": "2025-01-15T03:45:00Z",
  "from": "alice@example.com",
  "to": "bob@example.com",
  "reason": "shift change"
}
\end{minted}

\paragraph{Automatic transfer.}
Automatic transfer occurs when:
\begin{itemize}
  \item An \texttt{ops.incident} step is executed and the incident-commander takes over as DRI.
  \item An escalation timeout fires and the backup DRI is activated.
\end{itemize}

\paragraph{Same-actor-cannot-approve invariant.}
When the DRI transfers, the new DRI cannot approve their own gates. If alice was the DRI and bob takes over, bob cannot approve a gate that requires the \texttt{approver} role if bob is both the DRI and the approver. This separation-of-duties rule prevents self-approval.
